% Generated draft: requires mathematical and source review.
\documentclass[11pt]{article}
\usepackage[margin=25mm]{geometry}
\usepackage{fontspec}
\setmainfont{Latin Modern Roman}
\usepackage{amsmath,amssymb,mathtools,unicode-math}
\setmathfont{Latin Modern Math}
\usepackage{enumitem,needspace,xurl,hyperref}
\setlength{\parindent}{0pt}
\setlength{\parskip}{5pt}
\setlength{\emergencystretch}{3em}
\newcommand{\N}{\mathbb{N}}
\newcommand{\Z}{\mathbb{Z}}
\newcommand{\Q}{\mathbb{Q}}
\newcommand{\R}{\mathbb{R}}
\newcommand{\C}{\mathbb{C}}
\newcommand{\F}{\mathbb{F}}
\newcommand{\A}{\mathbb{A}}
\newcommand{\PP}{\mathbb{P}}
\newcommand{\E}{\mathbb{E}}
\newcommand{\eps}{\varepsilon}
\newcommand{\dd}{\,\mathrm{d}}
\begin{document}
\section*{UnsolvedMath ID: 800773}
\textbf{Source problem code:} catalogue entry
\subsection{Definitions and mathematical statement}
In the smooth complex projective interpretation suggested by the source title, let \(X\) be a rationally connected manifold of dimension \(n\). Here rational connectedness means that two general points of \(X\) lie on the image of a morphism \(\PP^1_{\C}\to X\). Write \(T_X\) for its tangent bundle and \(K_X\) for its canonical divisor, so that \(\mathcal O_X(-K_X)\cong\det T_X\). A divisor \(D\) on \(X\) is big if
\[
\limsup_{m\to\infty}\frac{n!\,\dim_{\C}H^0(X,\mathcal O_X(mD))}{m^n}>0.
\]
\textbf{Question (Matsumura, Problem 8, attributed there to [HIM19, Problem 3.13]).} If \(T_X\) is pseudo-effective, must \(-K_X\) be big? Equivalently, under the stated projective interpretation, does
\[
T_X\text{ pseudo-effective}
\quad\Longrightarrow\quad
\limsup_{m\to\infty}\frac{n!\,\dim_{\C}H^0(X,(\det T_X)^{\otimes m})}{m^n}>0
\]
hold for every rationally connected \(X\)?

\textbf{Unresolved qualifications.} The extracted question says only ``manifold''; smooth complex projectivity is the contextual interpretation used above, not an explicitly supplied hypothesis. The extract also gives no definition of pseudo-effectivity for vector bundles. That hypothesis is retained in the sense intended by [HIM19], whose definition must be checked before this becomes a fully specified standalone statement.

The supplied motivation attributes to [DPS94] the assertion that rationally connected manifolds with nef tangent bundle are Fano, meaning that their anticanonical divisor is ample. It says that the corresponding Fano conclusion fails for pseudo-effective tangent bundles and cites pseudo-effectivity of tangent bundles of Hirzebruch surfaces in [HIM19, Proposition 4.5]. These are source-reported contextual assertions. The preceding discussion mentions fibers \(F\) of a morphism \(\phi\colon X\to Y\), but the extracted question itself imposes no morphism or fiber hypothesis.
\subsection{Short English statement}
The question asks whether every rationally connected manifold with pseudo-effective tangent bundle has big anticanonical divisor. In the contextual smooth complex projective interpretation, this means that anticanonical sections have positive leading growth of order \(m^{\dim X}\). The ambient category and the intended definition of pseudo-effective tangent bundle require source verification.
\subsection{Sources}
\begin{itemize}
\item Local UnsolvedMath registry, supplied complete source record: permanent ID \(800773\), exact source code \texttt{catalogue entry}, title ``Positively Curved Projective Varieties --- Problem 8.'' The registry attributes the problem to Matsumura (2022) and labels it open; this is a supplied status, not an independently verified assessment.
\item Matsumura, \emph{Open Problems on Structure of Positively Curved Projective Varieties} (2022), Problem 8. Supplied source URL: \url{https://arxiv.org/abs/2211.09153}. The record reports access on 2026-10-16 and extraction by \texttt{arxiv\_latex}; these are supplied provenance details, not a claim of independent access.
\item [HIM19, Problem 3.13], identified in the supplied extract as the earlier source of the question. [HIM19, Proposition 4.5] is cited there for pseudo-effectivity of tangent bundles of Hirzebruch surfaces. Full bibliographic details and a separate URL are not supplied.
\item [DPS94], cited in the supplied extract for the Fano conclusion under rational connectedness and nefness of the tangent bundle. Full bibliographic details and a separate URL are not supplied.
\end{itemize}
\end{document}
